\documentclass[a4paper,12pt]{article}

\usepackage[ngerman]{babel}
\usepackage[utf8]{inputenc}
\usepackage[T1]{fontenc}
\usepackage{lmodern}
\usepackage{geometry}
\usepackage{listings}
\usepackage{xcolor}

\geometry{margin=2.5cm}

\title{Praktische Übung: Ein Legacy-System migrieren}
\author{Modul 321}
\date{\today}

\lstset{
    basicstyle=\ttfamily\small,
    breaklines=true,
    frame=single,
    backgroundcolor=\color{gray!10}
}

\begin{document}

\maketitle

\section*{Ziel der Übung}

In dieser Übung plant ihr eine \textbf{einfache Migration} eines bestehenden
Softwaresystems. Dabei wendet ihr die beiden Strategien aus der Lektion an:

\begin{itemize}
\item \textbf{Big Bang}
\item \textbf{Strangler Pattern}
\end{itemize}

Ihr sollt nicht ein komplettes Produkt bauen, sondern ein kleines Szenario
analysieren, eine Strategie wählen und sie mit einer einfachen technischen
Skizze umsetzen.

\section*{Ausgangslage}

Ein bestehendes Schulportal ist als ein älteres System aufgebaut.
Es enthält:

\begin{itemize}
\item Login
\item Kursübersicht
\item Nachrichten
\item Dateidownload
\end{itemize}

Die Kursübersicht soll modernisiert werden.
Der Rest darf vorerst im bestehenden System bleiben.

\section*{Projektidee}

Erstellt eine \textbf{kleine Demo-Architektur} für die Migration.
Die Demo kann als Zeichnung, HTML-Skizze oder kleines Node.js-Projekt umgesetzt werden.

\textbf{Minimalziel:}
\begin{itemize}
\item Ein Teil wird als \textbf{neu} markiert.
\item Ein Teil bleibt als \textbf{legacy}.
\item Ihr zeigt, wie Anfragen entweder direkt auf das alte System gehen
      oder auf einen neuen Teil umgeleitet werden.
\end{itemize}

\section*{Variante A: Skizze ohne Code}

Erstellt ein einfaches Diagramm oder ein Dokument mit:

\begin{itemize}
\item altem System
\item neuem Teil
\item Router / Gateway / Umschaltpunkt
\item Begründung der gewählten Strategie
\end{itemize}

\section*{Variante B: Kleine technische Demo}

Erstellt ein sehr kleines Projekt, zum Beispiel mit dieser Struktur:

\begin{lstlisting}
migration-demo/
|
|-- legacy/
|   `-- portal.js
|
|-- modern/
|   `-- course-overview.js
|
`-- gateway.js
\end{lstlisting}

\textbf{Beispielidee:}
\begin{itemize}
\item \texttt{legacy/portal.js} liefert alle alten Funktionen.
\item \texttt{modern/course-overview.js} liefert die neue Kursübersicht.
\item \texttt{gateway.js} entscheidet, wohin eine Anfrage geht.
\end{itemize}

\section*{Aufgaben}

\subsection*{Aufgabe 1: System verstehen}
Beantwortet kurz:
\begin{enumerate}
\item Welcher Teil ist \textbf{legacy}?
\item Welcher Teil wird zuerst ersetzt?
\item Warum ist genau dieser Teil ein sinnvoller Startpunkt?
\end{enumerate}

\subsection*{Aufgabe 2: Strategie wählen}
Entscheidet euch für eine Strategie:
\begin{itemize}
\item \textbf{Strangler Pattern} oder
\item \textbf{Big Bang}
\end{itemize}

Begründet die Wahl in \textbf{3 bis 5 Sätzen}.

\subsection*{Aufgabe 3: Umschaltlogik zeigen}
Zeigt praktisch oder als Pseudocode:

\begin{lstlisting}
if route == "/courses":
    use new module
else:
    use legacy system
\end{lstlisting}

Die Logik darf auch als Flussdiagramm dargestellt werden.

\subsection*{Aufgabe 4: Risiken benennen}
Nennen Sie mindestens:
\begin{itemize}
\item ein Risiko bei einem \textbf{Big Bang}
\item ein Risiko beim \textbf{Strangler Pattern}
\item eine Massnahme, wie dieses Risiko reduziert werden kann
\end{itemize}

\subsection*{Aufgabe 5: Nächster Migrationsschritt}
Beschreibt, welcher Teil nach der Kursübersicht als Nächstes migriert werden könnte
und warum.

\section*{Abgabe}

Die Abgabe soll enthalten:

\begin{itemize}
\item kurze Beschreibung des Ausgangssystems
\item gewählte Strategie mit Begründung
\item Skizze, Diagramm oder kleine Routing-Demo
\item kurze Reflexion zu Risiken und Reihenfolge
\end{itemize}

\section*{Reflexionsfragen}

\begin{enumerate}
\item Warum ist die Reihenfolge bei einer Migration wichtig?
\item Wann wäre ein Big Bang trotz Risiko sinnvoll?
\item Welche Vorteile bringt das schrittweise Vorgehen?
\item Was macht den Parallelbetrieb schwierig?
\end{enumerate}

\section*{Bewertung}

\begin{itemize}
\item Strategiewahl ist nachvollziehbar
\item Unterschied zwischen \textbf{legacy} und \textbf{neu} ist klar
\item Umschaltlogik ist verständlich dargestellt
\item Risiken und Gegenmassnahmen sind sinnvoll beschrieben
\item Die Reihenfolge der Migration ist begründet
\end{itemize}

\end{document}
